\section{性交后的身体护理与清洁}

性活动结束后的一小段自我照护，往往被忽略，却对预防感染、减少不适和保护长期性健康很有帮助。本节给出简明、可操作的事后护理要点。

\subsection{排尿与清洁}

\begin{itemize}
	\item \textbf{性交后尽快排尿}：这是预防"蜜月膀胱炎"（性交后尿路感染）的核心措施——排尿可冲走被推入尿道的细菌。女性尿道短而直，尤其容易发生上行感染。
	\item \textbf{清理分泌物}：用清水或温和的湿巾轻轻擦拭外阴与肛周即可，\textbf{不要冲洗阴道内部}（会破坏阴道菌群、增加感染风险）。
	\item \textbf{由前向后擦拭}：避免将肛周细菌带到尿道口与阴道口。
	\item \textbf{更换内裤}：若感到潮湿不适，可更换干爽内裤；避免穿紧身、不透气的衣物。
\end{itemize}

\subsection{身体与情绪}

\begin{itemize}
	\item \textbf{休息与补水}：性活动消耗体力，相当于一次中等强度运动；适当休息、补充水分有助于恢复。
	\item \textbf{肌肉与不适}：若出现酸胀或轻微不适，多可自行缓解；持续性疼痛、腹痛、异常出血应及时就诊。
	\item \textbf{事后安抚}：拥抱、依偎与简短交谈有助于情绪平稳回落，也有利于关系亲密（详见"性交后反应"相关章节）。
	\item \textbf{检查是否有破损或异物}：如发现避孕套破裂或滑脱，应及时评估补救避孕与性传播感染预防的需要。
\end{itemize}

\subsection{需要就医的信号}

\begin{itemize}
	\item 性交后反复出现尿频、尿急、尿痛，或伴发热、腰痛（提示尿路感染上行）。
	\item 性交后出血、持续疼痛、异常阴道分泌物或异味。
	\item 出现皮疹、瘙痒（可能为乳胶避孕套过敏、润滑剂刺激或精液过敏）。
\end{itemize}

\noindent\textit{【编者按】}良好的事后护理不是洁癖，而是把\textbf{清洁、观察、安抚}三件事做对，既保护身体，也维护关系。
